% GENERATED FILE. Edit canonical lessons, then run npm run generate:books.
% canonical-source-hash: fnv1a64:0b3fdf7b3164af9a
% canonical-lessons: BN-C172-kheladhula, BN-C172-chhatri, BN-C172-bord, BN-C172-chok, BN-C172-britti

\chapter{Sports, Female student, Board, Chalk, Scholarship}
\label{ch:bn-sports-female-student-board-chalk-scholarship}
\hlchaptermodality{172}

\begin{chapteropening}
\textbf{By the end of this chapter:} \emph{I can say sports, a female student, a board, chalk and a scholarship.}

Complete the last lesson of chapter 172: I can say sports, a female student, a board, chalk and a scholarship.
\end{chapteropening}

\section[\texorpdfstring{\bn{খেলাধুলা} (khelādhulā)}{খেলাধুলা (khelādhulā)}]{\bn{খেলাধুলা} (khelādhulā) — sports}

\label{lesson:BN-C172-kheladhula}



\begin{quote}
Before the new one: say the Bengali for paddy, then the Bengali for high tide.
\end{quote}

\subsection*{You'll want to know: \bn{খেলাধুলা}}
\textbf{\bn{খেলাধুলা}} — \emph{khelādhulā} — \textquotedblleft{}sports\textquotedblright{}.

\begin{cousinweb}
\textbf{\bn{খেলাধুলা}} (\emph{khelādhulā}) is \emph{khelā}, \textquotedblleft{}play\textquotedblright{}, which CDIAL lists under \emph{*khēll}, \textquotedblleft{}play\textquotedblright{}, one of a group of forms it reconstructs (CDIAL 3918), + \emph{dhulā}, \textquotedblleft{}dust\textquotedblright{}, under \emph{*dhūḍi}, \textquotedblleft{}dust\textquotedblright{}, beside Sanskrit \emph{dhūli} (CDIAL 6835): \textquotedblleft{}play and dust\textquotedblright{}.
\end{cousinweb}

\begin{grammarlens}[title={What you've built}]
One more word for this chapter.
\end{grammarlens}

\subsection*{Your turn}
\begin{itemize}
\raggedright
  \item \emph{Say it:} \emph{khelādhulā}
  \item \emph{Say it:} \emph{khelādhulā}, once more
  \item \emph{Say it:} say \emph{khelādhulā}
  \item \emph{Recall:} say the Bengali for paddy, then the Bengali for high tide, then say \emph{khelādhulā} again
\end{itemize}

\begin{tcolorbox}[breakable,colback=teal!4,colframe=teal!35!black,title={Before you move on}]
What does \bn{খেলাধুলা} mean? (\textquotedblleft{}Sports\textquotedblright{}.) Say it once more.
\end{tcolorbox}
\section[\texorpdfstring{\bn{ছাত্রী} (chhātrī)}{ছাত্রী (chhātrī)}]{\bn{ছাত্রী} (chhātrī) — a female student}

\label{lesson:BN-C172-chhatri}



\begin{quote}
Before the new one: say the Bengali for high tide, then the Bengali for sports.
\end{quote}

\subsection*{You'll want to know: \bn{ছাত্রী}}
\textbf{\bn{ছাত্রী}} — \emph{chhātrī} — \textquotedblleft{}a female student\textquotedblright{}.

\begin{cousinweb}
\textbf{\bn{ছাত্রী}} (\emph{chhātrī}) is Sanskrit \emph{chātrī}, \textquotedblleft{}female pupil\textquotedblright{}, taken over whole as a learned loan, a tatsama.
\end{cousinweb}

\begin{grammarlens}[title={What you've built}]
One more word for this chapter.
\end{grammarlens}

\subsection*{Your turn}
\begin{itemize}
\raggedright
  \item \emph{Say it:} \emph{chhātrī}
  \item \emph{Say it:} \emph{chhātrī}, once more
  \item \emph{Say it:} say \emph{chhātrī}
  \item \emph{Recall:} say the Bengali for high tide, then the Bengali for sports, then say \emph{chhātrī} again
\end{itemize}

\begin{tcolorbox}[breakable,colback=teal!4,colframe=teal!35!black,title={Before you move on}]
What does \bn{ছাত্রী} mean? (\textquotedblleft{}A female student\textquotedblright{}.) Say it once more.
\end{tcolorbox}
\section[\texorpdfstring{\bn{বোর্ড} (borḍ)}{বোর্ড (borḍ)}]{\bn{বোর্ড} (borḍ) — a board}

\label{lesson:BN-C172-bord}



\begin{quote}
Before the new one: say the Bengali for sports, then the Bengali for a female student.
\end{quote}

\subsection*{You'll want to know: \bn{বোর্ড}}
\textbf{\bn{বোর্ড}} — \emph{borḍ} — \textquotedblleft{}a board\textquotedblright{}.

\begin{cousinweb}
\textbf{\bn{বোর্ড}} (\emph{borḍ}) is a loan from English \emph{board}.
\end{cousinweb}

\begin{grammarlens}[title={What you've built}]
One more word for this chapter.
\end{grammarlens}

\subsection*{Your turn}
\begin{itemize}
\raggedright
  \item \emph{Say it:} \emph{borḍ}
  \item \emph{Say it:} \emph{borḍ}, once more
  \item \emph{Say it:} say \emph{borḍ}
  \item \emph{Recall:} say the Bengali for sports, then the Bengali for a female student, then say \emph{borḍ} again
\end{itemize}

\begin{tcolorbox}[breakable,colback=teal!4,colframe=teal!35!black,title={Before you move on}]
What does \bn{বোর্ড} mean? (\textquotedblleft{}A board\textquotedblright{}.) Say it once more.
\end{tcolorbox}
\section[\texorpdfstring{\bn{চক} (chôk)}{চক (chôk)}]{\bn{চক} (chôk) — chalk}

\label{lesson:BN-C172-chok}



\begin{quote}
Before the new one: say the Bengali for a female student, then the Bengali for a board.
\end{quote}

\subsection*{You'll want to know: \bn{চক}}
\textbf{\bn{চক}} — \emph{chôk} — \textquotedblleft{}chalk\textquotedblright{}.

\begin{cousinweb}
\textbf{\bn{চক}} (\emph{chôk}) is a loan from English \emph{chalk}.
\end{cousinweb}

\begin{grammarlens}[title={What you've built}]
One more word for this chapter.
\end{grammarlens}

\subsection*{Your turn}
\begin{itemize}
\raggedright
  \item \emph{Say it:} \emph{chôk}
  \item \emph{Say it:} \emph{chôk}, once more
  \item \emph{Say it:} say \emph{chôk}
  \item \emph{Recall:} say the Bengali for a female student, then the Bengali for a board, then say \emph{chôk} again
\end{itemize}

\begin{tcolorbox}[breakable,colback=teal!4,colframe=teal!35!black,title={Before you move on}]
What does \bn{চক} mean? (\textquotedblleft{}Chalk\textquotedblright{}.) Say it once more.
\end{tcolorbox}
\section[\texorpdfstring{\bn{বৃত্তি} (britti)}{বৃত্তি (britti)}]{\bn{বৃত্তি} (britti) — a scholarship}

\label{lesson:BN-C172-britti}



\begin{quote}
Before the new one: say the Bengali for a board, then the Bengali for chalk.
\end{quote}

\subsection*{You'll want to know: \bn{বৃত্তি}}
\textbf{\bn{বৃত্তি}} — \emph{britti} — \textquotedblleft{}a scholarship\textquotedblright{}.

\begin{cousinweb}
\textbf{\bn{বৃত্তি}} (\emph{britti}) is Sanskrit \emph{vṛtti}, \textquotedblleft{}livelihood, maintenance\textquotedblright{}, taken over whole as a learned loan, a tatsama.
\end{cousinweb}

\begin{grammarlens}[title={What you've built}]
One more word for this chapter.
\end{grammarlens}

\subsection*{Your turn}
\begin{itemize}
\raggedright
  \item \emph{Say it:} \emph{britti}
  \item \emph{Say it:} \emph{britti}, once more
  \item \emph{Say it:} say \emph{britti}
  \item \emph{Recall:} say the Bengali for a board, then the Bengali for chalk, then say \emph{britti} again
\end{itemize}

\begin{tcolorbox}[breakable,colback=teal!4,colframe=teal!35!black,title={Before you move on}]
What does \bn{বৃত্তি} mean? (\textquotedblleft{}A scholarship\textquotedblright{}.) Say it once more.
\end{tcolorbox}
